\documentclass[../main]{subfiles}
\usetikzlibrary{arrows.meta,decorations.markings,calc,patterns}
\begin{document}
\chapter{Sistema trifásico desbalanceado: cálculo de corrientes de fase y de neutro}
\section{Sistema trifásico desbalanceado}
\subsection{Introducción}
En un sistema trifásico \emph{estrella con neutro} desbalanceado, cada fase puede tener una impedancia distinta. Si se conocen las \emph{tensiones de fase} (línea–neutro) como fasores y las impedancias de carga por fase, las corrientes se obtienen con números complejos y la corriente de neutro resulta de la suma vectorial de las corrientes de fase.
\info{Idea clave}{Trabajamos con fasores complejos:

$$
\underline{I}_k=\frac{\underline{V}_{kN}}{\underline{Z}_k},\qquad k\in\{A,B,C\},\qquad 
\underline{I}_N=\underline{I}_A+\underline{I}_B+\underline{I}_C,\qquad
I_N=\left|\underline{I}_N\right|.
]
}

\subsection{Fórmulas fundamentales}

\begin{equation}
\underline{I}_A=\frac{\underline{V}_{AN}}{\underline{Z}_A},\qquad
\underline{I}_B=\frac{\underline{V}_{BN}}{\underline{Z}_B},\qquad
\underline{I}_C=\frac{\underline{V}_{CN}}{\underline{Z}_C}
\end{equation}

\begin{equation}
\underline{I}_N=\underline{I}_A+\underline{I}_B+\underline{I}_C,\qquad
I_N=\left|\underline{I}_N\right|
\end{equation}

\noindent
\textit{Nota:} \(\underline{V}=\!V\angle\varphi\) indica magnitud \(V\) y ángulo \(\varphi\) en grados. Las divisiones y sumas se realizan en forma compleja (rectangular o polar).

\subsection{Ejemplo}

\ejemplo{Ejemplo}{Dados
\[
\underline{V}_{AN}=220\angle 0^\circ\ \text{V},\quad 
\underline{V}_{BN}=215\angle -120^\circ\ \text{V},\quad
\underline{V}_{CN}=225\angle 120^\circ\ \text{V},
$$

$$
\underline{Z}_A=20+j10\ \Omega,\quad
\underline{Z}_B=18+j6\ \Omega,\quad
\underline{Z}_C=22+j11\ \Omega.
$$

Calcular $\underline{I}_A,\ \underline{I}_B,\ \underline{I}_C$ e $I_N$.

\textbf{Solución (esquema):}

$$
\underline{I}_A=\frac{220\angle 0^\circ}{\sqrt{20^2+10^2}\angle \arctan\!\left(\tfrac{10}{20}\right)}=
\frac{220\angle 0^\circ}{22.36\angle 26.565^\circ}=9.84\angle -26.565^\circ\ \text{A}.
$$

$$
\underline{I}_B=\frac{215\angle -120^\circ}{\sqrt{18^2+6^2}\angle \arctan\!\left(\tfrac{6}{18}\right)}=
\frac{215\angle -120^\circ}{18.97\angle 18.435^\circ}=11.33\angle (-138.435^\circ)\ \text{A}.
$$

$$
\underline{I}_C=\frac{225\angle 120^\circ}{\sqrt{22^2+11^2}\angle \arctan\!\left(\tfrac{11}{22}\right)}=
\frac{225\angle 120^\circ}{24.64\angle 26.565^\circ}=9.13\angle 93.435^\circ\ \text{A}.
$$

Sumar fasores: $\underline{I}_N=\underline{I}_A+\underline{I}_B+\underline{I}_C$ y finalmente $I_N=\lvert \underline{I}_N\rvert$.
}

\section{Representación gráfica}

\begin{center}
\begin{tikzpicture}
\draw\[-stealth] (0,0) -- (3,0) node\[right] {Fase A};
\draw\[-stealth] (0,0) -- (1.5,2.6) node\[above right] {Fase B};
\draw\[-stealth] (0,0) -- (-1.5,2.6) node\[above left] {Fase C};
\draw\[dashed,->] (0,0) -- (0,3) node\[above] {Neutro};
\node\[below] at (0,0) {Centro};
\end{tikzpicture}
\end{center}

\newpage

\actividad{Responder brevemente}
\begin{enumerate}
\item ¿Cómo se calculan $\underline{I}_A,\underline{I}_B,\underline{I}_C$ a partir de $\underline{V}_{kN}$ y $\underline{Z}_k$?
\item ¿Por qué $\underline{I}_N$ se obtiene por suma fasorial y no escalar?
\item ¿Qué efecto produce un desbalance de $\underline{Z}_k$ en $I_N$?
\item ¿Conviene operar en forma polar o rectangular? Justifique brevemente.
\end{enumerate}

\actividad{Resolver (usar números complejos)}
\textit{Tome secuencia $ABC$ y ángulos en grados. Indique $\underline{I}_A,\ \underline{I}_B,\ \underline{I}_C$ y $I_N$.}
\begin{enumerate}
\item $\underline{V}_{AN}=220\angle 0^\circ\ \text{V},\ \underline{V}_{BN}=228\angle -120^\circ\ \text{V},\ \underline{V}_{CN}=232\angle 120^\circ\ \text{V};\ \underline{Z}_A=18+j9\ \Omega,\ \underline{Z}_B=22+j8\ \Omega,\ \underline{Z}_C=20+j12\ \Omega.$

\item $\underline{V}_{AN}=230\angle 0^\circ\ \text{V},\ \underline{V}_{BN}=225\angle -120^\circ\ \text{V},\ \underline{V}_{CN}=220\angle 120^\circ\ \text{V};\ \underline{Z}_A=16+j12\ \Omega,\ \underline{Z}_B=24+j6\ \Omega,\ \underline{Z}_C=18+j9\ \Omega.$

\item $\underline{V}_{AN}=215\angle 0^\circ\ \text{V},\ \underline{V}_{BN}=215\angle -120^\circ\ \text{V},\ \underline{V}_{CN}=240\angle 120^\circ\ \text{V};\ \underline{Z}_A=25+j5\ \Omega,\ \underline{Z}_B=20+j10\ \Omega,\ \underline{Z}_C=15+j12\ \Omega.$

\item $\underline{V}_{AN}=225\angle 0^\circ\ \text{V},\ \underline{V}_{BN}=235\angle -118^\circ\ \text{V},\ \underline{V}_{CN}=230\angle 122^\circ\ \text{V};\ \underline{Z}_A=12+j8\ \Omega,\ \underline{Z}_B=28+j7\ \Omega,\ \underline{Z}_C=21+j9\ \Omega.$

\item $\underline{V}_{AN}=218\angle 0^\circ\ \text{V},\ \underline{V}_{BN}=222\angle -120^\circ\ \text{V},\ \underline{V}_{CN}=226\angle 120^\circ\ \text{V};\ \underline{Z}_A=30+j10\ \Omega,\ \underline{Z}_B=18+j12\ \Omega,\ \underline{Z}_C=26+j6\ \Omega.$

\item $\underline{V}_{AN}=240\angle 0^\circ\ \text{V},\ \underline{V}_{BN}=230\angle -120^\circ\ \text{V},\ \underline{V}_{CN}=235\angle 120^\circ\ \text{V};\ \underline{Z}_A=22+j11\ \Omega,\ \underline{Z}_B=16+j8\ \Omega,\ \underline{Z}_C=32+j10\ \Omega.$

\item $\underline{V}_{AN}=225\angle 0^\circ\ \text{V},\ \underline{V}_{BN}=225\angle -120^\circ\ \text{V},\ \underline{V}_{CN}=225\angle 120^\circ\ \text{V};\ \underline{Z}_A=18+j6\ \Omega,\ \underline{Z}_B=18+j9\ \Omega,\ \underline{Z}_C=18+j12\ \Omega.$

\item $\underline{V}_{AN}=210\angle 0^\circ\ \text{V},\ \underline{V}_{BN}=235\angle -120^\circ\ \text{V},\ \underline{V}_{CN}=220\angle 120^\circ\ \text{V};\ \underline{Z}_A=20+j4\ \Omega,\ \underline{Z}_B=25+j15\ \Omega,\ \underline{Z}_C=14+j7\ \Omega.$

\item $\underline{V}_{AN}=232\angle 0^\circ\ \text{V},\ \underline{V}_{BN}=228\angle -121^\circ\ \text{V},\ \underline{V}_{CN}=224\angle 119^\circ\ \text{V};\ \underline{Z}_A=27+j9\ \Omega,\ \underline{Z}_B=19+j5\ \Omega,\ \underline{Z}_C=23+j11\ \Omega.$

\item $\underline{V}_{AN}=220\angle 0^\circ\ \text{V},\ \underline{V}_{BN}=220\angle -120^\circ\ \text{V},\ \underline{V}_{CN}=235\angle 120^\circ\ \text{V};\ \underline{Z}_A=15+j5\ \Omega,\ \underline{Z}_B=15+j10\ \Omega,\ \underline{Z}_C=40+j8\ \Omega.$
\end{enumerate}

\end{document}
